\section{Experimentos}
\label{sec:experiments}

\subsection{Ambiente de Teste}

Todos os benchmarks foram executados em uma única máquina host com as especificações listadas
na Tabela~\ref{tab:hardware}. Os três contêineres de banco de dados e o cliente Java de
benchmark compartilharam essa máquina sem cargas de trabalho concorrentes durante as sessões
de teste.

\begin{table*}[t]
  \centering
  \caption{Hardware e versões de software do ambiente de teste.}
  \label{tab:hardware}
  \begin{tabular}{ll}
    \hline
    \textbf{Componente} & \textbf{Detalhes} \\
    \hline
    CPU       & Intel Core i5-11400F (6c/12t, 2,60--4,40\,GHz) \\
    RAM       & 32\,GB DDR4 \\
    Disco     & 500\,GB NVMe SSD \\
    SO        & NixOS (Linux 6.12) \\
    \hline
    IoTDB       & \texttt{apache/iotdb:1.3.0-standalone} \\
    InfluxDB    & \texttt{influxdb:2.7-alpine} \\
    TimescaleDB & \texttt{timescale/timescaledb:latest-pg15} \\
    Benchmark   & iot-benchmark~\cite{iotbenchmark} (Java, Maven) \\
    \hline
  \end{tabular}
\end{table*}

Os valores de utilização de CPU são expressos como porcentagem de um único núcleo lógico,
consistente com as convenções de reporte do \texttt{docker stats}. Um valor de 600\% indicaria
utilização total de todos os seis núcleos físicos; valores por banco de dados na faixa de
20--100\% indicam carga bem dentro do orçamento computacional disponível.

\subsection{Desempenho de Escrita}
\label{sec:write-results}

\subsubsection{Ingestão Sequencial de Linha de Base}

A Tabela~\ref{tab:write-baseline} mostra o desempenho de escrita de linha de base na escala
pequena sob a configuração ajustada. O IoTDB atinge 2,16\,M\,pts/s, aproximadamente 5$\times$
o InfluxDB e 12$\times$ o TimescaleDB.

\begin{table*}[t]
  \centering
  \caption{Desempenho de escrita de linha de base na escala pequena (ajustado). Vazão em pts/s.}
  \label{tab:write-baseline}
  \begin{tabular}{lrrrr}
    \hline
    \textbf{BD} & \textbf{pts/s} & \textbf{Lat.\,Média} & \textbf{Lat.\,P99} & \textbf{CPU\,Média} \\
    \hline
    IoTDB       & 2.163.809 &  1,73\,ms &   1,18\,ms &  5,1\% \\
    InfluxDB    &   397.894 & 12,00\,ms & 190,50\,ms &  8,5\% \\
    TimescaleDB &   174.787 & 27,70\,ms & 223,00\,ms & 22,5\% \\
    \hline
  \end{tabular}
\end{table*}

Três fatores arquiteturais cumulativos explicam a diferença.

\textbf{Profundidade do protocolo.} O protocolo SESSION\_BY\_TABLET do IoTDB submete um tablet
colunar tipado diretamente por uma conexão TCP persistente usando serialização binária Thrift.
Uma syscall submete um lote completo diretamente no buffer de escrita do TsFile. O InfluxDB
requer uma requisição HTTP por escrita, incluindo gerenciamento de conexão, análise de texto do
protocolo de linha e pesquisa de índice TSI. O TimescaleDB roteia cada escrita pela pilha
completa do PostgreSQL: TCP, análise SQL, planejamento de consulta, roteamento de hipertabela,
escrita no WAL e formatação de tupla MVCC. Cada camada adiciona latência; a pilha PostgreSQL
adiciona mais.

\textbf{Buffering em memória e compromissos de durabilidade.} O IoTDB acumula tablets em seu
heap JVM e reconhece escritas antes da persistência em disco, descarregando para TsFile
assincronamente quando o memtable se enche. Isso troca explicitamente a garantia de durabilidade
por vazão, um compromisso que deve ser avaliado por implantação. O cache TSM do InfluxDB
opera de forma similar, mas é limitado (8\,GB na configuração ajustada). O TimescaleDB segue
a semântica padrão do PostgreSQL: cada escrita é descarregada pelo WAL antes do
reconhecimento, fornecendo a garantia de durabilidade mais forte ao custo de E/S síncrona
por escrita.

\textbf{Manutenção de índices.} O caminho de escrita do IoTDB é thin acima de seu buffer JVM.
O motor TSM do InfluxDB mantém um Índice de Séries Temporais (TSI) a cada escrita. Os índices
B-tree e BRIN do TimescaleDB crescem com o volume de dados, e a atividade de autovacuum
aumenta a sobrecarga de escrita à medida que a hipertabela se expande.

A combinação desses três fatores indica que a vantagem de escrita do IoTDB tem base
arquitetural: o protocolo binário, o buffering assíncrono e a ausência de manutenção
de índice síncrona atuam de forma cumulativa, não sendo eliminados por ajuste de
configuração.

\begin{figure}[ht]
  \centering
  \includegraphics[width=\columnwidth]{figures/pt_br/chart_write_throughput_small}
  \caption{Vazão de escrita (pts/s, escala logarítmica) para todas as quatro cargas de trabalho
  de escrita na escala pequena (ajustado). O IoTDB lidera em todos os tipos de teste.}
  \label{fig:write_throughput_small}
\end{figure}

\subsubsection{Sensibilidade ao Tamanho do Lote}

A Tabela~\ref{tab:batch} mostra como a vazão muda entre BATCH\_SIZE=1 (uma requisição de
escrita por ponto) e BATCH\_SIZE=1000 (1.000 pontos por requisição).

\begin{table}[ht]
  \centering
  \caption{Sensibilidade ao tamanho do lote na escala pequena (ajustado). Vazão em pts/s.}
  \label{tab:batch}
  \begin{tabular}{lrrr}
    \hline
    \textbf{BD} & \textbf{LOTE=1} & \textbf{LOTE=1000} & \textbf{Razão} \\
    \hline
    IoTDB       &  28.593 & 11.827.572 & 413$\times$ \\
    InfluxDB    &   5.024 &  2.667.256 & 531$\times$ \\
    TimescaleDB &   5.120 &    214.927 &  42$\times$ \\
    \hline
  \end{tabular}
\end{table}

IoTDB e InfluxDB mostram aceleração de 413--531$\times$, refletindo o alto custo da sobrecarga
por requisição. Para o InfluxDB, trata-se de uma rodada HTTP mais sobrecarga de conexão TCP;
para o IoTDB, é o custo fixo de uma chamada Thrift por tablet. A razão de 42$\times$ do
TimescaleDB é marcadamente menor porque o flush do WAL do PostgreSQL já domina no LOTE=1:
a sobrecarga por linha é tão alta que lotes maiores fornecem benefício proporcionalmente menor;
o gargalo passa da sobrecarga de protocolo para a E/S de armazenamento.

Este resultado tem significado prático direto. Uma aplicação IoT que envia uma leitura de
sensor por requisição HTTP ao InfluxDB atingirá vazão 531$\times$ menor do que aquela que
armazena localmente em buffer e escreve em lotes. Em qualquer volume de escrita onde o lote
seja viável, é a escolha de configuração de maior impacto para usuários de InfluxDB e IoTDB.

\subsubsection{Ingestão Fora de Ordem}

A Tabela~\ref{tab:ooo} mostra como cada banco de dados lida com a carga de trabalho de 20\%
de dados com chegada atrasada.

\begin{table*}[t]
  \centering
  \caption{Penalidade fora de ordem: vazão (pts/s) e degradação relativa em relação à escrita
  de linha de base, escala pequena (ajustado).}
  \label{tab:ooo}
  \begin{tabular}{lrrr}
    \hline
    \textbf{BD} & \textbf{ESCRITA\,(pts/s)} & \textbf{FORA-ORDEM\,(pts/s)} & \textbf{Penalidade} \\
    \hline
    IoTDB       & 2.163.809 & 1.871.351 & $-$13,5\% \\
    InfluxDB    &   397.894 &   255.634 & $-$35,7\% \\
    TimescaleDB &   174.787 &   160.121 & $-$8,4\% \\
    \hline
  \end{tabular}
\end{table*}

O TimescaleDB mostra a menor degradação relativa porque as atualizações de índice B-tree do
PostgreSQL são independentes da ordem: inserir uma linha com um instante de tempo anterior
incorre no mesmo custo que inserir uma na fronteira de escrita atual. A penalidade de 13,5\%
do IoTDB surge porque sua árvore LSM deve compactar entradas desordenadas para restaurar a
ordem classificada exigida pelos blocos TsFile. O InfluxDB incorre na maior penalidade
relativa (35,7\%) por causa de regravações de blocos TSM necessárias para acomodar pontos
fora de sequência.

Em vazão absoluta, o IoTDB ainda lidera amplamente sob condições fora de ordem
(1,87\,M vs.\ 255\,K pts/s para InfluxDB). Portanto, duas dimensões devem ser consideradas
separadamente: vazão absoluta fora de ordem (IoTDB lidera em cada escala) e degradação
relativa por desordem de escrita (TimescaleDB é o mais resiliente com a menor penalidade
fracional).

\begin{figure}[ht]
  \centering
  \includegraphics[width=\columnwidth]{figures/pt_br/chart_ooo_degradation}
  \caption{Degradação de vazão fora de ordem (\%) relativa à escrita sequencial de linha de
    base, em todas as três escalas. O B-tree do TimescaleDB permanece estável; as
    regravações de blocos TSM do InfluxDB consistentemente incorrem na maior penalidade
  relativa.}
  \label{fig:ooo_degradation}
\end{figure}

\subsection{Desempenho de Leitura}
\label{sec:read-results}

Uma observação importante sobre a interpretação da vazão de leitura: o valor em pts/s para um
teste de leitura soma as contagens de pontos retornados por todos os tipos de consulta ativos.
Uma consulta de intervalo que retorna centenas de pontos por chamada parecerá ter vazão em pts/s
muito maior do que uma consulta de último valor que retorna exatamente um ponto por chamada.
Esses valores não são, portanto, comparáveis entre tipos de teste de leitura.
A \textbf{latência} é usada exclusivamente como a métrica primária de comparação de leitura
nesta seção.

Uma observação metodológica relevante para todas as comparações de leitura: os testes de leitura
são executados sequencialmente no mesmo conjunto de dados dentro de uma sessão, de modo que
testes posteriores se beneficiam de caches de páginas do SO mais quentes e estruturas de banco
de dados em memória (ver Seção~\ref{sec:limitations}). Onde os efeitos de cache quente são o
mecanismo declarado para o desempenho de um banco de dados (como o \texttt{shared\_buffers}
do TimescaleDB para subamostragem), isso é uma propriedade esperada e reproduzível sob
condições normais de operação, não um artefato.

\subsubsection{Consultas de Último Valor}

A Tabela~\ref{tab:latest} mostra a latência média de último valor em todas as três escalas
para as execuções ajustadas.

\begin{table}[ht]
  \centering
  \caption{Latência média de último valor (ms) por escala (ajustado). Menor é melhor.}
  \label{tab:latest}
  \begin{tabular}{lrrr}
    \hline
    \textbf{BD} & \textbf{Pequena} & \textbf{Média} & \textbf{Grande} \\
    \hline
    TimescaleDB &  0,38 &  0,41 &  0,62 \\
    IoTDB       &  1,24 &  0,51 &  0,68 \\
    InfluxDB    &  9,24 & 21,72 & 23,24 \\
    \hline
  \end{tabular}
\end{table}

TimescaleDB e IoTDB mantêm latência abaixo de um milissegundo ou próxima disso em todo o
crescimento de cem vezes no volume de dados. Ambos exploram atalhos específicos do banco de dados
para recuperação do último valor. O TimescaleDB executa um \texttt{SELECT~...~ORDER~BY~time~DESC~LIMIT~1}
contra o bloco de hipertabela mais recente via varredura B-tree reversa; como os blocos de
hipertabela são ordenados por tempo, isso sempre toca a menor quantidade possível de dados.
O IoTDB mantém um cache de último valor em memória dedicado por série, servindo a consulta
inteiramente da RAM sem acessar o TsFile.

O InfluxDB não tem atalho equivalente. Ele deve buscar o final dos blocos de arquivo TSM de
cada série, e esse custo de busca cresce à medida que mais arquivos TSM se acumulam ao longo
do tempo de vida do conjunto de dados. A progressão resultante (9,24\,ms na escala pequena,
21,72\,ms na média, 23,24\,ms na grande) confirma uma relação aproximadamente linear entre
dados armazenados e custo de último valor. Na janela de escalas avaliada, a latência de último valor do InfluxDB mostrou crescimento
aproximadamente linear com o volume de dados armazenados. Em implantações sem política
ativa de retenção ou downsampling, esse padrão pode torná-lo desvantajoso para
monitoramento em tempo real de alta frequência; implantações que apliquem essas políticas
podem mitigar esse comportamento.

\begin{figure}[ht]
  \centering
  \includegraphics[width=\columnwidth]{figures/pt_br/chart_latest_point_scaling}
  \caption{Latência média de último valor (ms) por escala. TimescaleDB e IoTDB permanecem
    próximos de zero (caminhos de acesso O(1)); o InfluxDB cresce linearmente à
  medida que a contagem de arquivos TSM aumenta.}
  \label{fig:latest_point_scaling}
\end{figure}

\subsubsection{Subamostragem}

A Tabela~\ref{tab:downsample} mostra a latência de subamostragem por janela de tempo com
\texttt{GROUP~BY}.

\begin{table}[ht]
  \centering
  \caption{Latência média de subamostragem (ms) por escala (ajustado). Menor é melhor.}
  \label{tab:downsample}
  \begin{tabular}{lrrr}
    \hline
    \textbf{BD} & \textbf{Pequena} & \textbf{Média} & \textbf{Grande} \\
    \hline
    TimescaleDB & 0,61 & 0,52 & 0,73 \\
    IoTDB       & 1,81 & 0,64 & 0,95 \\
    InfluxDB    & 5,71 & 5,65 & 9,22 \\
    \hline
  \end{tabular}
\end{table}

A agregação \texttt{time\_bucket()} do TimescaleDB permanece abaixo de 1\,ms em todo o
crescimento de cem vezes no volume de dados, resultado que se repete de forma consistente em todas as escalas testadas. A
chave é que as consultas \texttt{time\_bucket()} varrem uma janela de tempo fixa em vez do
conjunto de dados completo: uma vez que os blocos relevantes estão quentes em
\texttt{shared\_buffers}, o custo da consulta é limitado pelo tamanho da janela, não pelo
total de dados armazenados. Essa propriedade torna o TimescaleDB excelente para cargas de
trabalho de painel que consistentemente consultam janelas de tempo recentes independentemente
da acumulação de dados históricos.

A latência de subamostragem do IoTDB melhora da escala pequena para a média
(1,81\,ms $\to$ 0,64\,ms) à medida que a execução mais longa na escala média dá ao compilador
JIT da JVM tempo para otimizar o caminho quente de agregação de blocos. Na escala grande,
permanece abaixo de 1\,ms. O motor de agregação Flux do InfluxDB incorre em maior sobrecarga
por consulta (5--9\,ms) por custos de avaliação de expressões e de pipeline de operadores.

\subsubsection{Consultas de Filtro de Valor}

Consultas de filtro de valor (predicados sobre valores de sensores como
\texttt{temperature > 85.0}) são a carga de trabalho mais discriminatória arquiteturalmente
neste estudo. Índices baseados em tempo não podem respondê-las diretamente; cada motor deve
varrer todos os valores armazenados ou explorar uma estrutura de dados secundária para eliminar
blocos de armazenamento que não correspondem.

A Tabela~\ref{tab:vfilter} mostra a latência de filtro de valor e seu comportamento de
escalonamento por escala.

\begin{table*}[t]
  \centering
  \caption{Latência média de filtro de valor (ms) por escala (ajustado).
    A coluna de razão mostra o crescimento de latência da escala pequena para a grande
  (dados cem vezes maiores).}
  \label{tab:vfilter}
  \begin{tabular}{lrrrr}
    \hline
    \textbf{BD} & \textbf{Pequena} & \textbf{Média} & \textbf{Grande} & \textbf{Razão P$\to$G} \\
    \hline
    IoTDB       &   3,72 &    9,80 &   19,42 &  5,2$\times$ \\
    TimescaleDB &  23,98 &  130,08 &  510,05 & 21,3$\times$ \\
    InfluxDB    &  71,09 &  310,75 &  496,81 &  7,0$\times$ \\
    \hline
  \end{tabular}
\end{table*}

A vantagem do IoTDB é grande e \textit{melhora} com a escala. Cada bloco TsFile armazena
estatísticas de mínimo/máximo por série. Um predicado de valor pode pular blocos inteiros cujo
intervalo armazenado não intercepta o limiar do predicado, fornecendo indexação secundária
leve na camada de armazenamento. Na escala grande, com mais dados por bloco, uma fração maior
de blocos pode ser eliminada, produzindo a razão de escalonamento sublinear (5,2$\times$ para
dados cem vezes maiores). Esse é o comportamento esperado de uma estrutura de salto baseada
em estatísticas: a razão de eliminação melhora à medida que os blocos ficam mais densos.

O TimescaleDB usa o BRIN (Block Range INdex) do PostgreSQL, que registra valores de
mínimo/máximo por intervalo de bloco em disco. O BRIN fornece eliminação com granularidade
mais grosseira do que as estatísticas por bloco do IoTDB e degrada em taxa superlinear
(21,3$\times$ para dados cem vezes maiores). Na escala grande, o TimescaleDB fica em média 510\,ms
por consulta de filtro de valor, uma desvantagem de 26$\times$ relativa ao IoTDB. Essa
carga de trabalho também leva o TimescaleDB à sua maior utilização de CPU (77\% médio,
83\% de pico), indicando que o executor de consultas é o gargalo.

O InfluxDB não tem estrutura de eliminação baseada em estatísticas, mas se beneficia de algum
salto grosseiro no nível de arquivo baseado em metadados de intervalo de tempo no índice TSM,
produzindo uma razão de escalonamento de 7$\times$, significativamente melhor do que a do
TimescaleDB, mas ainda 10$\times$ pior do que a do IoTDB em termos absolutos na escala grande.

\begin{figure}[ht]
  \centering
  \includegraphics[width=\columnwidth]{figures/pt_br/chart_value_filter_scaling}
  \caption{Latência média de filtro de valor (ms, escala logarítmica) por escala. As
    estatísticas de mínimo/máximo por bloco do IoTDB produzem crescimento sublinear
    (5,2$\times$ para dados cem vezes maiores); o índice BRIN do TimescaleDB degrada
  superlinearmente (21,3$\times$).}
  \label{fig:value_filter_scaling}
\end{figure}

\subsubsection{Consultas de Intervalo}

As consultas de intervalo recuperam todos os pontos de dados dentro de uma janela de tempo
sem filtragem por valor. O TimescaleDB é o mais rápido aqui: 0,54\,ms na escala pequena e
0,66\,ms na grande, beneficiando-se da eliminação de blocos de hipertabela (apenas blocos que
sobrepõem a janela de consulta são abertos) e do cache de páginas do SO quente que evita E/S
em disco para janelas de tempo recentes frequentemente consultadas. O IoTDB atinge
1,19--1,21\,ms e é efetivamente insensível à escala dentro deste intervalo de dados. O
InfluxDB atinge 5,41--8,68\,ms, onde o scanner de blocos TSM suporta maior sobrecarga por
arquivo do que a varredura sequencial no nível de bloco do TimescaleDB.

\subsection{Footprint de Memória}
\label{sec:memory-results}

As medições de memória requerem interpretação cuidadosa porque as execuções originais do
benchmark com início compartilhado produziram valores substancialmente inflados por um artefato
metodológico.

\textbf{TimescaleDB} reportou aproximadamente 5,4\,GB desde o primeiro teste nas execuções
originais. Isso não era uso do pool de buffers do PostgreSQL; era cache de páginas do SO
acumulado das sessões de escrita anteriores do IoTDB. O Linux contabiliza o cache de páginas
do sistema de arquivos como parte do RSS de um contêiner no \texttt{docker stats}, e o
contêiner ocioso do TimescaleDB estava em execução há horas enquanto o IoTDB gerava gigabytes
de E/S. Nas execuções sequenciais ajustadas, o TimescaleDB começa em 239\,MB e cresce para
aproximadamente 1,8\,GB após os testes de escrita aqueceram o \texttt{shared\_buffers}.
O valor de 1,8\,GB é o footprint operacional preciso para esse tamanho de conjunto de dados
sob a configuração ajustada.

\textbf{IoTDB} reportou aproximadamente 12,5\,GB fixo ao longo das execuções originais.
Durante as horas em que ficou ocioso enquanto InfluxDB e TimescaleDB eram testados, sua JVM
pré-alocou especulativamente o heap máximo configurado em antecipação a carga futura. Nas
execuções ajustadas, a JVM começa em 1,7\,GB e cresce organicamente para aproximadamente
10\,GB ao final da fase de leitura, refletindo a demanda real do conjunto de trabalho por
fase de carga em vez de reserva especulativa.

\textbf{InfluxDB} é o mais eficiente em memória entre as cargas de trabalho de escrita:
190--940\,MB em testes de escrita ajustados na escala pequena. No entanto, essa vantagem não
se mantém incondicionalmente. Na escala grande, o LOTE-GRANDE atinge pico de 2,9\,GB à medida
que o cache de escrita TSM mantém grandes volumes de lotes pendentes junto com buffers de
compactação concorrente. A afirmação comumente citada de que a memória do InfluxDB é ``limitada
a 1\,GB por padrão'' descreve apenas o teto do cache TSM; o RSS total do processo inclui o
índice de séries, buffers WAL e sobrecarga do runtime Go que juntos podem exceder
substancialmente esse limite sob alta carga de escrita.

A implicação prática para o planejamento de implantação é que a vazão de escrita dominante do
IoTDB requer um compromisso substancial de memória ($\sim$10\,GB na escala média, crescendo
para $\sim$15\,GB na grande). Para ambientes com memória restrita (gateways de borda IoT,
sistemas embarcados ou instâncias de nuvem com RAM limitada), o footprint mais enxuto do
InfluxDB pode ser o fator decisivo mesmo quando a vazão de escrita é a principal preocupação.

\begin{figure}[ht]
  \centering
  \includegraphics[width=\columnwidth]{figures/pt_br/chart_memory_scaling}
  \caption{Footprint de memória de pico (MB) por escala (execuções ajustadas). O conjunto de
    trabalho JVM do IoTDB cresce para 15\,GB na escala grande. O InfluxDB permanece a
  opção mais eficiente em memória em cada escala.}
  \label{fig:memory_scaling}
\end{figure}

\subsection{Efeitos de Escala}
\label{sec:scale-effects}

\begin{figure*}[t]
  \centering
  \includegraphics[width=\textwidth]{figures/pt_br/chart_heatmap_latency}
  \caption{Mapa de calor de degradação de latência por tipos de carga de trabalho de leitura e
    transições de escala (execuções ajustadas). Verde indica latência estável; vermelho
    indica alto crescimento. O IoTDB é consistentemente verde exceto para filtro de valor
    (pequena$\to$média); o TimescaleDB mostra degradação severa no filtro de valor em
  ambas as transições.}
  \label{fig:heatmap_latency}
\end{figure*}

\subsubsection{Escalonamento da Vazão de Escrita}

A Tabela~\ref{tab:write-scale} mostra como a vazão de escrita de linha de base evolui por
escalas nas execuções ajustadas.

\begin{table}[ht]
  \centering
  \caption{Vazão de escrita de linha de base (pts/s) por escala (ajustado).}
  \label{tab:write-scale}
  \begin{tabular}{lrrr}
    \hline
    \textbf{BD} & \textbf{Pequena} & \textbf{Média} & \textbf{Grande} \\
    \hline
    IoTDB       & 2.163.809 &  7.991.224 & 23.607.475 \\
    InfluxDB    &   397.894 &    392.975 &  1.256.602 \\
    TimescaleDB &   174.787 &    175.685 &    400.455 \\
    \hline
  \end{tabular}
\end{table}

A vazão de escrita do IoTDB escala superlinearmente: 2,2\,M pts/s na escala pequena,
8,0\,M na média, 23,6\,M na grande. Isso é impulsionado pela compilação JIT da JVM:
benchmarks curtos não dão ao JIT tempo suficiente para otimizar totalmente todos os caminhos
quentes. Nas escalas média e grande, as execuções mais longas permitem que o JIT atinja a
compilação de pico, saturando o caminho de escrita Thrift. Os números de escrita do IoTDB na
escala pequena, portanto, subestimam seu teto em estado estacionário; o valor na escala grande
é o mais representativo da vazão de produção sustentada.

O InfluxDB permanece aproximadamente constante entre as escalas pequena e média (WAL e
compactação TSM são o gargalo independentemente do tamanho do conjunto de dados), mas melhora
marcadamente na escala grande à medida que volumes de escrita maiores permitem que o motor de
compactação TSM amortize sua sobrecarga de forma mais eficiente. O TimescaleDB mostra um
padrão similar: plano da escala pequena para a média, depois melhorando na escala grande à
medida que a execução mais longa permite mais intervalos de checkpoint e lotes individuais
maiores para se acumularem.

\begin{figure}[ht]
  \centering
  \includegraphics[width=\columnwidth]{figures/pt_br/chart_write_scaling}
  \caption{Vazão de escrita de linha de base (pts/s, escala logarítmica) por escala (ajustado).
    O crescimento superlinear do IoTDB reflete a saturação JIT da JVM; InfluxDB e
  TimescaleDB permanecem planos até a escala grande.}
  \label{fig:write_scaling}
\end{figure}

\subsubsection{Saturação JIT do IoTDB no LOTE-PEQUENO}

A vazão de \texttt{batch-small} do IoTDB por escalas (28\,K $\to$ 81\,K $\to$ 446\,K pts/s)
é uma das demonstrações mais claras de escalonamento impulsionado por JIT no conjunto de dados.
A latência média simultaneamente \textit{diminui} de 1,22\,ms para 0,40\,ms à medida que a
escala aumenta, o que só é possível se a melhoria vier de otimização de código em vez de
redução de carga. O valor na escala grande (446\,K pts/s, 0,40\,ms de latência média)
representa a vazão de ingestão de ponto único em estado estacionário do IoTDB uma vez que o
JIT compilou completamente o caminho quente de serialização Thrift.

\subsubsection{Timeout do LOTE-GRANDE do TimescaleDB}

O teste \texttt{batch-large} do TimescaleDB atinge o timeout de 3.600 segundos do benchmark
na escala grande tanto nas execuções original quanto nas ajustadas. Com 50 dispositivos e um
tamanho de lote de 1.000 linhas, cada requisição de escrita toca 50 partições de hipertabela
simultaneamente e dispara até 50 flushes WAL concorrentes. O custo cumulativo de E/S esgota
o orçamento de tempo antes que a contagem total de laços seja completada. Este não é um caso
extremo patológico, mas uma consequência previsível do fan-out de partição de hipertabela sob
escritas de lote grande com alta concorrência. Para cargas de trabalho de carregamento em lote
direcionadas ao TimescaleDB em escala, a contagem de partições deve ser ponderada contra a
amplificação de flush WAL.

\subsubsection{Crescimento Linear de Último Valor do InfluxDB}

Entre as tendências multiescala, a progressão de latência de último valor do InfluxDB
(9,24 $\to$ 21,72 $\to$ 23,24\,ms) merece atenção particular porque revela um problema
estrutural de escalabilidade. Cada arquivo TSM incremental escrito em disco adiciona à
profundidade de busca que uma consulta de último valor deve percorrer. TimescaleDB e IoTDB
mostram latência de último valor plana ou até \textit{decrescente} (devido ao aquecimento JIT
para IoTDB) ao longo do mesmo crescimento de cem vezes no volume de dados. Qualquer implantação que
acumulará meses ou anos de dados verá o custo de último valor do InfluxDB continuar crescendo,
enquanto os custos do TimescaleDB e do IoTDB permanecem limitados.

\subsection{Impacto do Ajuste de Configuração}
\label{sec:tuning-results}

\subsubsection{Ganhos do TimescaleDB com pgtune}

A Tabela~\ref{tab:tsdb-tuning} compara o desempenho de escrita do TimescaleDB antes e depois
do pgtune na escala pequena.

\begin{table*}[t]
  \centering
  \caption{Desempenho de escrita do TimescaleDB: original vs.\ ajustado (escala pequena).
  $\Delta$ mostra variação percentual; negativo significa melhoria para latência.}
  \label{tab:tsdb-tuning}
  \begin{tabular}{lrrrr}
    \hline
    \textbf{Teste} & \textbf{pts/s ant.} & \textbf{pts/s dep.} & $\Delta$vazão & $\Delta$lat \\
    \hline
    ESCRITA        & 133.677 & 174.787 & $+$30\% & $-$25\% \\
    LOTE-GRANDE    & 160.546 & 214.927 & $+$34\% & $-$25\% \\
    FORA-DE-ORDEM  & 132.189 & 160.121 & $+$21\% & $-$21\% \\
    FILTRO-VALOR   & (latência) & (latência) & ---     & $-$40\% \\
    LEITURA        & (latência) & (latência) & ---     & $-$38\% \\
    \hline
  \end{tabular}
\end{table*}

O principal driver é \texttt{shared\_buffers=8GB}, que fornece ao PostgreSQL um pool de
buffers grande o suficiente para evitar evicções de buffer que forçam escritas extras no WAL.
O \texttt{work\_mem=27413kB} elimina derramamentos em disco para avaliação de predicados
baseada em ordenação, explicando a redução de 40\% na latência de filtro de valor, uma
carga de trabalho que requer ordenação de resultados filtrados. O
\texttt{checkpoint\_completion\_target=0.9} distribui a E/S de checkpoint de forma mais
uniforme ao longo do intervalo de checkpoint, reduzindo os picos de paralisação de escrita
que inflam a latência P99 (o P99 do LOTE-GRANDE caiu de 1.069\,ms para 812\,ms). Essas
melhorias indicam que as configurações padrão do PostgreSQL são intencionalmente conservadoras
para garantir funcionamento em hardware de baixo recurso, e que o ajuste via pgtune ou
ferramenta equivalente deve ser aplicado antes de reportar resultados em hardware de médio e
alto porte.

\subsubsection{Ganhos do InfluxDB com Isolamento}

O InfluxDB não recebeu alterações de configuração nas execuções ajustadas. Suas melhorias
vêm inteiramente da execução sem bancos de dados concorrentes. Todas as latências de leitura
melhoraram 15--21\%, consistente com melhor utilização do cache de arquivos TSM uma vez que
o cache de páginas do TimescaleDB não compete mais pelo espaço de buffer do SO. O resultado
mais dramático foi o P99 do LOTE-GRANDE: 205\,ms $\to$ 60\,ms ($-$71\%). Na configuração
original, o TimescaleDB estava escrevendo agressivamente em disco durante a sessão de teste
do InfluxDB, privando periodicamente o motor de compactação TSM do InfluxDB de largura de
banda de E/S. Essas paralisações raras mas severas de E/S produziram outliers de latência
extremos que não afetaram a vazão média, mas inflaram o P99 dramaticamente. O isolamento
removeu essas paralisações.

\subsubsection{Correção de Dois Falsos Positivos na Escala Grande}

As execuções originais na escala grande produziram duas descobertas que as reexecuções
ajustadas reverteram.

\textbf{Penalidade fora de ordem do TimescaleDB.} A execução original na escala grande
mostrou uma penalidade de 23\% na vazão fora de ordem para o TimescaleDB, atribuída a ciclos
de descompressão-recompressão de inserções fora de ordem em blocos comprimidos mais antigos.
A reexecução ajustada mostra apenas 2,7\% de degradação, consistente com o padrão nas
escalas pequena e média. O mecanismo: inserções fora de ordem em blocos mais antigos requerem
a leitura desses blocos do disco para descomprimí-los. Na sessão original com início
compartilhado, o heap JVM em expansão do IoTDB (mais de 15\,GB) havia despejado as páginas
do cache do SO que continha esses blocos, forçando leituras frias do disco para cada inserção
fora de ordem em um bloco mais antigo. Com acesso exclusivo, os blocos permanecem quentes no
cache e os custos de descompressão são negligíveis.

\textbf{Saturação de CPU do TimescaleDB.} A carga de trabalho READ original na escala grande
apresentou média de 97,7\% de utilização de CPU com latência P99 de 4.232\,ms, levando à
conclusão de que o TimescaleDB ``ficou sem capacidade computacional'' e não poderia ser
melhorado sem uma CPU mais forte ou fragmentação horizontal. A reexecução ajustada mostra
77,4\% de CPU médio e P99 de 1.850\,ms. O valor original de 97,7\% foi o resultado de todos
os três bancos de dados e o cliente Java de benchmark competindo simultaneamente pelos mesmos
seis núcleos de CPU. O TimescaleDB não estava em um teto de hardware: estava privado de CPU
por processos concorrentes.

Essas duas reversões ressaltam a importância de uma metodologia de benchmark isolada e
corretamente ajustada. Sem ciclos de vida sequenciais de contêineres e configuração padrão de
banco de dados, ambas as afirmações teriam sido publicadas como descobertas validadas.

\begin{figure*}[t]
  \centering
  \includegraphics[width=\textwidth]{figures/pt_br/chart_read_saturation_comparison}
  \caption{Utilização média de CPU da carga de trabalho READ (acima) e latência P99 (abaixo,
    escala logarítmica) para execuções original (hachurado) vs.\ ajustado (sólido) em
    todas as três escalas. A CPU na escala grande do TimescaleDB cai de 98\% para 77\%
    e o P99 de 4.233\,ms para 1.850\,ms uma vez que os contêineres concorrentes são
  removidos.}
  \label{fig:read_saturation_comparison}
\end{figure*}

\subsection{Carga de Trabalho Mista de Leitura}

A carga de trabalho \texttt{read} executa todos os dez tipos de consulta do iot-benchmark com
peso igual, fornecendo uma linha de base de leitura composta que inclui consultas de intervalo,
agregações, subamostragem, consultas de último valor e consultas de predicado de valor
simultaneamente. Na escala pequena (ajustado), o TimescaleDB tem média de 0,67\,ms, o IoTDB
de 2,11\,ms e o InfluxDB de 14,28\,ms por operação no conjunto combinado. Esses valores são
consistentes com os resultados individuais de consulta apresentados acima: o TimescaleDB
domina para consultas de janela de tempo, que constituem a maioria dos dez tipos, enquanto a
vantagem do IoTDB nas consultas de filtro de valor é diluída pela ponderação igual.

O papel analítico primário da carga de trabalho mista neste estudo é como benchmark de ajuste
(Tabela~\ref{tab:tsdb-tuning}): a latência READ do TimescaleDB melhorou 38\% após o pgtune,
confirmando que múltiplos tipos de consulta (não apenas o filtro de valor) se beneficiam
de um pool de buffers dimensionado corretamente.

\subsection{Resumo}
\label{sec:results-summary}

A Tabela~\ref{tab:summary} consolida o mapeamento de vencedor por carga de trabalho em todos
os cenários testados.

\begin{table*}[t]
  \centering
  \caption{Resumo de recomendação por carga de trabalho em todos os cenários testados.}
  \footnotesize
  \label{tab:summary}
  \begin{tabular}{l l p{0.55\linewidth}}
    \hline
    \textbf{Carga de Trabalho} & \textbf{Melhor BD} & \textbf{Razão arquitetural principal} \\
    \hline
    Ingestão sequencial (ordenada)          & IoTDB       & Protocolo de tablet binário + buffering em memória antes do flush para TsFile \\
    Carregamento em lote (LOTE=1000)        & IoTDB       & Caminho de escrita colunar vetorizado; 11,8\,M pts/s na escala pequena \\
    Ingestão de ponto único (LOTE=1)        & IoTDB       & Menor sobrecarga por chamada Thrift vs. rodada HTTP \\
    Fora de ordem (vazão absoluta)          & IoTDB       & 1,87\,M pts/s OOO vs. 255\,K (InfluxDB); lidera em cada escala \\
    Fora de ordem (tolerância relativa)     & TimescaleDB & Atualizações de B-tree são independentes de ordem \\
    Subamostragem em painel                 & TimescaleDB & \texttt{time\_bucket()} em \texttt{shared\_buffers} quente; latência constante \\
    Exportação de intervalo de tempo        & TimescaleDB & Eliminação de blocos + cache de páginas quente do SO \\
    Alertas por limite (filtro de valor)    & IoTDB       & Estatísticas de mín/máx por bloco; escalonamento sublinear com volume de dados \\
    Consulta de último valor em tempo real  & TimescaleDB & Varredura B-tree reversa sub-ms no bloco mais recente \\
    Eficiência de memória (escrita)         & InfluxDB    & Cache TSM limitado; sem pré-alocação JVM \\
    \hline
  \end{tabular}
\end{table*}

O IoTDB é a escolha clara para cargas de trabalho dominadas por escrita, mas requer compromisso
substancial de memória ($\sim$10\,GB na escala média, crescendo para $\sim$15\,GB na grande)
e não suporta junções relacionais. O TimescaleDB troca velocidade bruta de escrita por
compatibilidade com SQL, o ecossistema completo do PostgreSQL e desempenho de leitura
consistentemente forte para consultas de janela de tempo uma vez que os dados estão
armazenados em cache. O InfluxDB ocupa o meio-termo: melhor vazão de escrita do que o
TimescaleDB, melhor eficiência de memória do que o IoTDB, mas mais fraco em consultas
complexas de predicado de valor e consultas de último valor em escala. Nenhum banco de dados
único domina em todas as cargas de trabalho.
